\section{System Model and Measurements}
\label{sec:model}

We study unicast application messages in a packet-level LoRa mesh model.
Each flow $i$ has unique identity, source $s_i$, destination $d_i$,
enqueue time $a_i$, and relative deadline $\Delta_i$:
\begin{equation}
D_i=a_i+\Delta_i.
\label{eq:deadline}
\end{equation}
The source observes its local flow state, cached path, radio queue,
configured PHY and timers, and decoded packets. Relays additionally use
received SNR and their configured role to schedule forwarding. No protocol
receives future faults, global reachability, or a destination-delivery
oracle. Offline graph diagnostics characterize the environment only.

Let $t_i^{D}$ be the first eligible destination reception and $t_i^{A}$
the first accepted source ACK; missing observations have time infinity.
For all $N$ scheduled flows, timely delivery and confirmation are
\begin{equation}
P_D=\frac{1}{N}\sum_i\mathbf{1}[t_i^{D}\le D_i],\qquad
P_A=\frac{1}{N}\sum_i\mathbf{1}[t_i^{A}\le D_i].
\label{eq:outcomes}
\end{equation}
Capacity rejection, no-start flows, and deadline misses remain in $N$.
An overheard routed packet is eligible delivery only when the destination
is its designated next hop. Delivery without timely confirmation is an
outcome, not by itself a diagnosis of ACK-channel loss.

For a fixed generation-plus-drain cutoff $C$, charge every distinct physical
transmission $q$ starting at $t_q\le C$:
\begin{equation}
A=\sum_{q:t_q\le C}T(L_q),\qquad
C_D=\frac{A}{\sum_i\mathbf{1}[t_i^{D}\le D_i]}.
\label{eq:physical-cost}
\end{equation}
$T(L)$ is full-frame LoRa time on air, using the standard symbol and
payload-length calculation~\cite{augustin2016study}. All DATA, ACK,
discovery, duplicate, and recovery transmissions are charged. A frame
starting before the cutoff is charged in full even if it ends later.
$C_D$ is undefined with zero timely deliveries. Requests and cancellations
are separate events; neither is a physical transmission.

For CPR, a path carrying $n$ node identifiers, useful payload of $P$ bytes,
and optional stage marker $m$ has modeled lengths
\begin{equation}
\begin{aligned}
L_{D/F}(n,m)&=12+2n+P+\mathbf{1}[m\ne\varnothing],\\
L_A(n,m)&=12+2n+\mathbf{1}[m\ne\varnothing].
\end{aligned}
\label{eq:wire}
\end{equation}
The ten-byte common header plus path count/index fields and two-byte node
IDs define this custom encoding. CPR's request ID equals its flow ID and
adds no discovery-ID extension. Cache generations and confirmation revisions
are local state. These lengths do not reproduce official product framing,
encryption, or authentication overhead.

\section{MeshEcho-CPR}
\label{sec:cpr}

% Reproducible vector diagram. Logical order and schematic timing only;
% positions are not measured latency or drawn to time scale.
\begin{figure*}[t]
\centering
\begin{tikzpicture}[
  x=1cm,y=1cm,>=Latex,
  stage/.style={draw,rounded corners=1pt,align=center,font=\small,
    text width=2.35cm,minimum height=0.78cm,inner sep=4pt},
  annotation/.style={align=center,font=\footnotesize},
  flow/.style={->,line width=0.65pt}]

\node[anchor=west,font=\bfseries\small] at (0,3.5)
  {(a) Source-stage order for full CPR};
\node[stage] (select) at (1.25,2.1) {Initial local\\opportunity};
\node[stage] (initial) at (5,2.1) {Routed initial\\marker absent, $t_0$};
\node[stage] (repeat) at (9,2.1) {Same-path repeat\\marker 1, $t_1$};
\node[stage] (recovery) at (13,2.1) {Recovery flood\\marker 2};
\draw[flow] (select) -- node[annotation,above,yshift=12pt] {READY route\\and budget fit} (initial);
\draw[flow] (initial) -- node[annotation,above,yshift=12pt] {No ACK yet\\observe at $t_0+B_1$} (repeat);
\draw[flow] (repeat) -- node[annotation,above,yshift=12pt] {No ACK yet\\observe at $t_1+B_1$} (recovery);

\node[stage] (cold) at (5,0.55) {Initial DATA flood\\marker absent};
\draw[flow] (select.south) |- (cold.west);
\node[annotation] at (2.5,1.08) {Otherwise,\\if flood fits};
\node[annotation,anchor=west,text width=6.7cm] at (6.8,0.55)
  {No repeat or recovery on this branch.\\
   Each later stage on the routed branch rechecks route and deadline at its proposed start.};
\node[annotation,text width=15cm] at (7.1,-0.55)
  {A valid ACK from any actually started eligible stage completes the live flow.
   Deadline expiry also retires it.\\
   Both stop pending source work; committed sends remain charged.
   Rejected admission creates no physical TX.};

\node[anchor=west,font=\bfseries\small] at (0,-1.4)
  {(b) One relay copy: commitment can precede physical start};
\draw[->,line width=0.65pt] (0.5,-2.25) -- (14.5,-2.25) node[right,font=\small] {time};
\fill[black!12] (2,-2.1) rectangle (6,-2.4);
\fill[black!38] (6,-2.1) rectangle (10,-2.4);
\fill[black!75] (10,-2.1) rectangle (13,-2.4);
\foreach \x/\lab in {2/{request $r_q$},6/{commit $k_q$},10/{start $t_q$},13/{end $e_q$}} {
  \draw (\x,-2.06) -- (\x,-2.53);
  \node[annotation,below] at (\x,-2.55) {\lab};
}
\node[annotation] at (4,-1.85) {Uncommitted};
\node[annotation] at (8,-1.85) {Reserved in queue};
\node[annotation] at (11.5,-1.85) {Physical TX};
\draw[->,dashed] (4,-3.65) -- (4,-2.42);
\node[annotation,below,text width=4.2cm] at (4,-3.65)
  {Matching off-path ACK can cancel this pending copy};
\draw[->,dashed] (8,-3.65) -- (8,-2.42);
\node[annotation,below,text width=4.8cm] at (8.7,-3.65)
  {ACK is too late to cancel, even before physical start};
\end{tikzpicture}
\caption{Logical source stages and the relay commitment boundary. In (a),
$B_1=B_R(p_i,1)$; observation times are capped at $D_i$, and a new stage
cannot start at $D_i$. Arrows show conditional stage order, not guaranteed
transmissions. The no-repeat control reserves the nominal observation slot
before recovery. In (b), the time axis is schematic, with
$t_q=\max(k_q,b_v)$ and $e_q=t_q+T(L_q)$; no latency measurements are
encoded. Eligibility also requires the matching conditions in
Eq.~\eqref{eq:cancel}.}
\label{fig:controller-timing}
\end{figure*}


\subsection{Local admission and finite source stages}

A queued send has request-creation time, commitment time $k_q$, and reserved
physical start $t_q$. With the local radio reserved until $b_v$,
\begin{equation}
t_q=\max(k_q,b_v),\qquad e_q=t_q+T(L_q).
\label{eq:commit}
\end{equation}
An admitted request commits at $k_q$ and reserves $[t_q,e_q)$ immediately.
It cannot be canceled later, even when $t_q>k_q$. Admission hooks receive
the proposed start, so queue delay enters the source's deadline check.

Let $h=|p|-1$ be cached-path hops, $H$ the hop limit, and $\epsilon$ a
safety margin. The source computes a route budget
\begin{equation}
B_R(p,m)=h\{T(L_D(h+1,m))+T(L_A(h+1,m))\}+\epsilon.
\label{eq:route-budget}
\end{equation}
For flood base delay $b$ and uniform jitter width $j$, it uses
\begin{equation}
\begin{aligned}
B_F(m)={}&H\{T(L_F(H+1,m))+T(L_A(H+1,m))\}\\
&+(H-1)(b+j+0.35)+\epsilon.
\end{aligned}
\label{eq:flood-budget}
\end{equation}
The 0.35-s term is the configured model's maximum weak-link delay addition.
The synthetic maximum path sizes this admission budget; actual transmitted
frames are charged at their actual path lengths. An action fits when
\begin{equation}
\operatorname{fits}(t,B,D)\iff t<D\ \land\ t+B\le D.
\label{eq:fits}
\end{equation}
These are local admission estimates. They do not bound all relay queues,
interference, loss, or successful ACK return time.

At the initial opportunity, the source chooses \texttt{R\_INITIAL} only
when its cached path is READY, unexpired, and feasible with the unmarked
route budget. Otherwise it attempts \texttt{F\_INITIAL} using the unmarked
flood budget. This initial-flood branch includes absent, quarantined,
expired, and locally infeasible routes. If neither action fits, no frame
starts and the flow remains in the denominator.

For a routed initial start $t_0$, define $B_1=B_R(p_i,1)$. At
$\min(D_i,t_0+B_1)$, a still-live, unconfirmed flow can request one repeat
on exactly the saved path. At its proposed start $t_1$, the repeat requires
the same current READY route generation, unexpired state, and a feasible
marked-frame budget. Only after this repeat has a recorded physical start
does full CPR observe for recovery at $\min(D_i,t_1+B_1)$. A still-unconfirmed
flow may then request one recovery flood, again requiring the same READY
generation and its own proposed-start budget. Failed repeat admission does
not skip directly to recovery. An initial-flood flow arms neither stage.
Figure~\ref{fig:controller-timing} separates this source-stage order from
the relay's request, commitment, and physical-start times.

Diagnostic action names alone do not identify rejection causes. In the
implemented recovery branch, the action
\texttt{DROP\_DEADLINE\_INFEASIBLE} is also emitted with reason
\texttt{stale-route-generation} when the saved route is no longer current
and READY, even if the remaining deadline admits the flood budget. Such
an event rejects that recovery start; it does not demonstrate a shortage
of deadline slack or immediately erase the live flow. Cause-specific counts
must inspect the reason and corresponding route/time state.

The initial packet is unmarked, the repeat has marker 1, and recovery has
marker 2. All retain the same flow ID and deadline. A valid ACK from any
actually started eligible stage completes the flow once and cancels its
remaining uncommitted source requests. Deadline expiry also retires pending
work; committed transmissions remain charged. Equal-time decoded receptions
precede deadline timers, permitting an eligible ACK at $D_i$, while a new
source stage cannot start at $D_i$.

If $N_i(m)$ counts DATA/FLOOD transmissions by the original source,
\begin{equation}
\begin{gathered}
N_i(m)\le1,\quad m\in\{\varnothing,1,2\},\\
\sum_mN_i(m)\le
\begin{cases}1,&\text{initial flood},\\3,&\text{initial route}.
\end{cases}
\end{gathered}
\label{eq:source-bound}
\end{equation}
Unique flow identity, once-only stage guards, and rejection of a flood
returning to its source establish this bound. It excludes ACKs and other
nodes' forwarding; it is not a network-wide three-transmission bound.

\subsection{Forwarding and message confirmation}

Routed DATA follows the designated next hop. A FLOOD is processed only if
its flow/request IDs agree, its marker is supported, and its simple path
starts at the origin and ends at the decoded sender with the correct path
index. The receiver must be absent from the path and cannot be the origin.
For incoming path length $n$, require $1\le n\le H$ and
$1\le\mathrm{TTL}\le H-n+1$. Reception appends the receiver; a destination
can ACK on the last permitted hop, whereas further forwarding requires
remaining TTL and relay eligibility. Initial and recovery duplicate keys
remain distinct. Duplicate suppression cancels only an uncommitted handle.

An ACK traverses the accepted path in reverse. At the source, it must match
a live flow, reversed endpoints, flow/request identity, supported stage,
simple path, source index, and decoded preceding transmitter; it must carry
no application payload and arrive by $D_i$. Its stage must have a recorded
start no later than the ACK reception. Routed initial/repeat ACKs must
match the saved path exactly; initial-flood/recovery ACKs may confirm their
corresponding discovered path. Replay after retirement cannot complete a
flow again. These are honest-packet simulator checks, not authentication.

Message confirmation is separate from shared-cache mutation. Store a route
as $(p,g,\rho,\mathrm{expiry},\mathrm{phase})$, with generation $g$ and
confirmation revision $\rho$. A valid ACK for the current READY, unexpired
path refreshes expiry, preserves $g$, and increments $\rho$. Otherwise its
path may replace the cache only if the stage's captured destination epoch
still matches the current epoch. A stale ACK can therefore confirm its
message without overwriting a newer different route. Shared generation
changes are not ACK-rejection conditions. A routed-flow timeout may
quarantine a READY route only if both saved $g$ and $\rho$ still match;
an older timeout cannot invalidate a newer same-path confirmation, including
one recorded at the same timestamp.

\subsection{Relay cancellation and component controls}

Each pending relay record holds endpoints, flow/request IDs, marker,
observed path, and the exact send handle. A locally decoded ACK $a$ can
cancel handle $q$ at relay $v$ only when
\begin{equation}
\begin{aligned}
\operatorname{cancel}(v,q,a)={}&\operatorname{decoded}(v,a)
\land\operatorname{match}(q,a)\\
&\land[v\notin\operatorname{path}(a)]\\
&\land\operatorname{uncommitted}(q)
\land\neg\operatorname{canceled}(q).
\end{aligned}
\label{eq:cancel}
\end{equation}
Matching includes endpoints, identity, marker, empty payload, a simple
correctly oriented ACK path, valid reverse index, and decoded sender.
The relay must be off the ACK path. It need not know the source's current
cache generation or deadline. Count decoded ACK checks separately from
successful cancellations: checks can encounter no pending handle or a
committed one. Their ratio is not success among eligible pending copies,
and canceled-copy counts are not network-wide counterfactual airtime savings.

The ACK-only control disables relay ACK cancellation while retaining ledger
capacity, eviction, source ACK processing, and duplicate suppression. The
no-repeat control emits no repeat and reserves the nominal observation slot,
opening recovery eligibility at $t_0+2B_R(p_i,1)$; recovery still passes its
own deadline and route checks. Queue changes can make later timelines
differ. The no-recovery control retains the initial transmission and repeat.
The historical whole-ledger-off control changes a broader policy and is
not an ACK-only ablation. Clean controls require resolved parameters to
differ only in their named switch.

The authoritative route and active-flow capacities are eight per source;
the cancellation ledger holds eight records per relay and evicts its oldest
entry, canceling it only if uncommitted. Separate inherited limits remain:
128 pending flood entries and 128 seen keys per node, with a 30-s seen TTL.
Destination-epoch storage is bounded by route plus active-flow capacity;
the retired-flow ring stores 256 IDs per protocol instance. These component
limits do not bound the entire simulator's event/evidence memory. Path
validation costs $O(H)$; the current shared Python ledger scans can cost
$O(N_{\mathrm{nodes}}K)$ for ledger capacity $K$, so whole operations are
not claimed constant-time.
